%% kapa: exam=1 round=36 year=2025 period=1 form=A subject=civil_law subject_ko=민법 seq=12 number=12
%% source: Q-Net 감정평가사 기출문제 (cst00309) artlSeq=5231525 file=raw/1008/5231525/2025년 제36회 감정평가사 1차시험 1교시 문제지 원본.pdf
%% conversion: kordoc 4.24.51 -> tools/gen1.py (auto); review=proofread-sonnet; flags=none
\begin{question}{12}
甲은 자신 소유의 X토지에 관하여 乙과 유효하게 매매계약(제1 계약)을 체결한 후 乙로부터 매매대금 전액을 지급받았다. 丙은 제1 계약과 관련한 甲의 배임행위에 적극 가담하여 甲과 X토지에 관한 매매계약(제2 계약)을 체결한 후, 이를 원인으로 甲으로부터 소유권이전등기를 경료받았다. 이에 관한 설명으로 옳은 것을 모두 고른 것은? (다툼이 있으면 판례에 따름)
\begin{framedtext}
ㄱ. X토지에 관하여 경료된 丙명의의 소유권이전등기는 원인무효이다. \par
ㄴ. 만약 丙이 X토지를 丁에게 매도하고 소유권이전등기를 경료하여 주었다면, 丁은 제2 계약이 유효하다고 주장할 수 있다. \par
ㄷ. 만약 丙이 X토지를 무단점유하고 있는 戊에 대하여 X토지의 소유권에 기한 반환을 주장하는 경우, 戊는 제2 계약의 무효를 주장할 수 있다.
\end{framedtext}
\begin{choices}
  \item ㄱ
  \item ㄴ
  \item ㄱ, ㄷ
  \item ㄴ, ㄷ
  \item ㄱ, ㄴ, ㄷ
\end{choices}
\end{question}
\answer{3}
